\begin{thesisabstract}{english}

Large language models running in multi-step workflows pay for every
token they read: energy, latency, and cost scale with context
length, and much of that context carries little information. Context
compression is the natural lever, and the single-agent
prompt-compression literature is mature. What has not been measured
is whether a compressor that preserves single-agent
question-answering accuracy also preserves the harder property a
planner needs --- combining information distributed across several
compressed fragments into a correct answer. This thesis measures
that on \emph{multi-fragment LLM workflows}, using a deterministic
regex solver or a single LLM call over all compressed fragments as
the planner; multi-round agent simulation is outside the empirical
scope.

The headline contribution is the H1 finding above. Three further
contributions support it. \textbf{C1} is a reproducible
$150$-instance multi-fragment benchmark across three workload
families (cross-document fact aggregation, constraint-satisfaction
planning, multi-step retrieval). \textbf{C2} is the empirical
characterisation itself --- a sharp \emph{coordination cliff}
$\tau^{*}$ across four training-free compressors, together with a
compounding-error model whose contribution is explaining \emph{why}
the transition is sharp (a recall threshold on critical tokens
produces a step in success); position prediction is a
\emph{first-order} estimate only (median relative error $\approx
35\%$, predicted band misses the empirical $\tau^{*}$ on $11/11$
testable cells). A direct A3 probe with an LLM planner --- the only
non-circular test of the threshold mechanism --- confirms the step
shape on the dense family ($k\approx 15$, $r_0 \approx 0.84$) and a
graded transition on the distributed family ($k \approx 5.9$, band
$\approx 0.54$). \textbf{C3} is a three-pipeline RAG catalogue with
a first-pass EUR-per-workflow cost ledger whose corpus-dominated
fidelity limits the P3-vs-P1 comparison (no cost-parity claim is
made). \textbf{C4} is a memory-bus reference implementation
together with a summary-level \emph{inference-disclosure} metric;
the metric is the C4 \emph{research} contribution and the bus is
the engineering artefact that operationalises it. CAAC operationalises
the compounding-error bound as a runtime back-off and serves as a
constructive demonstration of the cliff measurement rather than a
fifth headline contribution.

The central result is \textbf{H1}: at the workload level ($n=150$),
the change in single-agent QA-F1 under compression does not
positively predict the change in coordination success for any of the
four compressors (Spearman $\rho \in [-0.82, +0.32]$; all
$95\%$~bootstrap CIs exclude $0.6$). The single-agent benchmarks
that dominate the compression literature therefore mis-rank
compressors for multi-fragment use. A sharp cliff exists on $11$ of
$12$ (compressor, family) cells. Within a defined calibrated regime
the \emph{LLM-planner} cliff shows no detected shift across a single
well-identified scale-and-architecture comparison (local
Llama-3.1-8B $\tau^{*}\approx 2.48$ vs frontier Qwen-2.5-72B
$\approx 2.79$ on one cell); a third frontier model (DeepSeek~V4~Pro)
corroborates only weakly with a bootstrap CI too wide to identify
the position. The RAG sign-flip hypothesis is not supported ---
compress-first is preferred over retrieve-first on the single
retriever/embedder stack tested --- and compression reduces
inference disclosure, but in proportion to how aggressively it
destroys source tokens rather than through privacy-specific
filtering (an oracle field-redaction control achieves the same
$\sim 23$~pp reduction at near-zero compression, matching the
truncation baseline at $-24.6$~pp). The H1 result on family-a should
be read with one structural caveat: that family's regex solver
succeeds iff enough critical numeric tokens survive compression, and
the critical-token-recall metric \emph{is} that surviving fraction,
so the model's central equation on family-a relates two near-identical
measurements --- it is most directly \emph{computable} on family-a,
not most independently \emph{validating}.

In the interest of honesty: of the six pre-registered hypotheses,
two hold as stated (H1, H2) and the rest were reframed or not
supported as originally written (the H3 sign-flip; the H5
monotonicity claim, reframed as ceiling--cliff separation; the H6
transfer claim, reframed as information-density scaling); the
compounding-error model is a first-order, not a tight, position
predictor. These reframings are reported explicitly. The deliverable
is the manuscript and an open-source reference implementation: the
local pipeline reproduces from one command given a precomputed
compression cache, while the frontier-API arm is reproducible in
distribution rather than byte-for-byte.

\end{thesisabstract}
